\section{Introduction}

LLM-guided fuzzers are often described as bringing their own knowledge: a
model trained on large amounts of code is expected to know what inputs
stress a parser or a compiler~\cite{deng2023titanfuzz,xia2024fuzz4all}. A
previous study of differential testing for Dart JSON
deserialization~\dartcite{} found the opposite in a controlled setting. A
bounded refinement loop, in which an LLM writes and rewrites a
Hypothesis~\cite{maciver2019hypothesis} generator from a score-only
feedback signal, found cross-implementation divergence in only 2.9\% of
generated documents, against 22.2\% for a hand-built static generator.
Richer feedback, incremental mutation and a larger model did not close the
gap. Four observations taken from a thirteen-case manual characterization
did: with them in the prompt, the same loop reached 28.8\%. But it never
found a divergence it had not been told about. The bottleneck was domain
knowledge, and it came from the researchers.

This paper asks where else that knowledge can come from. A human
characterization is slow and requires an expert; if a fuzzer is to be
applied to a new ecosystem, the knowledge has to be obtained some other way.
Two automated sources are natural. An agent can \emph{read the code} of the
implementations under test and summarize what it finds, as white-box LLM
testing tools do~\cite{yang2024whitefox,yang2025kernelgpt,vadayath2026aijon}.
Or an agent can \emph{run experiments}: design a small number of probe
documents, observe how each implementation handles them, and write down what
it learned. The second is what the researchers did by hand in that study, and it
needs no source code, which matters when the implementations are closed,
generated at compile time, or reflective. We call it \emph{black-box
self-characterization}.

We compare these sources under control. The refinement loop, model,
sampling settings and feedback are those of that study, unchanged;
the only thing that differs between arms is the text in a single knowledge
slot of the prompt, capped at 1{,}500 characters. The four arms are
\emph{none}, \emph{human} (the four observations, verbatim),
\emph{code} (a summary written after reading the four paths'
deserialization code) and \emph{probe} (a summary written after running 20
probe documents in two rounds). The design, seeds, tests and decision rules
were fixed and committed before any code or API spend, and every later
change is recorded in a deviations log with its reason. We ask:

\begin{itemize}
\item \textbf{RQ1 (acquisition).} Can a probing agent acquire knowledge
that closes the refinement gap as well as human knowledge does? We answer
this on Dart, where the answer key is known.
\item \textbf{RQ2 (sources).} How do the four knowledge sources compare in
divergence rate, in which known divergence patterns they lead to, and in
cost?
\item \textbf{RQ3 (transfer and discovery).} On an ecosystem with no human
characterization, Kotlin with four widely used JSON libraries, does probing
knowledge help, does it rediscover a known behavior, and does it find
anything new?
\item \textbf{RQ4 (remedies, pre-registered extension).} If unguided
probing fails, is it because the probes miss the relevant inputs or because
the model is too weak? We test a probing agent told to cover generic test
categories first, and the unchanged agent with a larger model.
\end{itemize}

Our findings are:

\begin{itemize}
\item On Dart, knowledge sources are clearly ordered: human (30.9\%) above
code-read (11.6\%) above none (1.5\%), each difference significant after
correction. Probing reaches 10.5\% on average but is not distinguishable
from no knowledge ($p \approx 0.69$). It is all-or-nothing: 49\% on the one
seed whose probes hit a divergence, at most 1.3\% on the others.
\item The divergence patterns each arm finds follow the facts its summary
states. Every code summary states that two paths accept a non-integer
\texttt{id}, and every code run finds the \texttt{id} pattern; no code
summary states that one path substitutes an empty list for a missing
\texttt{tags} field, and one code run in five finds that pattern. No probe summary mentions non-integer \texttt{id}
values, and no probe run finds the \texttt{id} pattern.
\item On Kotlin, where 11--14 of 20 probes diverge on every seed, probing
knowledge helps: 67.1\% against 50.0\% over ten seeds ($p \approx 0.007$),
using a replication declared before it ran.
\item Both remedies help on Dart. Systematic probing reaches 20.2\%
($p \approx 0.032$ against none; $0.16$ after Holm correction across the
six declared tests), and every one of its runs finds the missing-field
pattern; a larger model reaches 12.1\%. Neither reaches the human arm's mean.
The agents covered the first checklist items thoroughly and never reached
the 64-bit boundary items, so the integer-overflow pattern stayed mostly
unfound: coverage decides what is learned, and a budget of 20 probes does not
cover enough.
\item Triage of every Kotlin divergence class found documented defaults and
known issues, including the Gson null-safety bypass that served as our
answer key, and one standards violation in kotlinx.serialization's strict
mode, which we reported upstream.
\end{itemize}

The overall lesson is conditional. Black-box self-characterization works
when the behaviors worth knowing are easy to stumble on. When they are
narrow, which is where an unguided fuzzer also fails and knowledge matters
most, it works only as far as its probes are made to reach them. The data,
agents, prompts and analysis are available at \url{\artifacturl}.
